\documentclass{exam-zh}
\usepackage{edu-practice}

\newtcolorbox{solutionblock}{
  enhanced, breakable,
  colback=edu-blue-bg,
  colframe=edu-blue!25,
  boxrule=0pt,
  borderline west={2pt}{0pt}{edu-blue!50},
  arc=0pt,
  left=3mm, right=3mm, top=2mm, bottom=2mm,
  before skip=2mm, after skip=3mm,
  fontupper=\small,
}
\newcounter{solstep}
\newcommand{\solstep}[2]{%
  \stepcounter{solstep}%
  \par\smallskip
  \noindent\hangindent=6mm
  {\color{edu-blue}\bfseries\textcircled{\scriptsize\thesolstep}\enspace #1}\enspace #2\par
}

\begin{document}
\section*{三角形一边平行线练习 · 答案与解析}
\section*{练习 1 答案}
\subsection*{第 1 题}
\setcounter{solstep}{0}
\begin{solutionblock}
\textbf{答案：}$AB=4$。\par\medskip
    \solstep{先判断}{已知边与未知边中出现平行线段 $CD$、$AB$，先求两个相似三角形的比。}
    \solstep{计算并约分}{$\dfrac{PA}{PC}=\dfrac{3}{6}=\dfrac{1}{2}$。}
    \solstep{标份数}{约分后，$AB$ 对应 1 份，$CD$ 对应 2 份。}
    \solstep{按份数公式求边}{$AB=CD\times\dfrac{1}{2}=8\times\dfrac{1}{2}=4$。}
\end{solutionblock}
\subsection*{第 2 题}
\setcounter{solstep}{0}
\begin{solutionblock}
\textbf{答案：}$BP=4$。\par\medskip
    \solstep{先判断}{已知边与未知边都不在平行线段 $AB,CD$ 上，求两侧对应线段的比。}
    \solstep{计算并约分}{$\dfrac{AP}{PC}=\dfrac{4}{6}=\dfrac{2}{3}$。}
    \solstep{标份数}{约分后，$BP$ 对应 2 份，$PD$ 对应 3 份。}
    \solstep{按份数公式求边}{$BP=PD\times\dfrac{2}{3}=6\times\dfrac{2}{3}=4$。}
\end{solutionblock}
\subsection*{第 3 题}
\setcounter{solstep}{0}
\begin{solutionblock}
\textbf{答案：}$PA=6$。\par\medskip
    \solstep{先判断}{已知边与未知边都不在平行线段 $AB,CD$ 上，求两侧对应线段的比。}
    \solstep{计算并约分}{$\dfrac{PB}{PD}=\dfrac{6}{8}=\dfrac{3}{4}$。}
    \solstep{标份数}{约分后，$PA$ 对应 3 份，$PC$ 对应 4 份。}
    \solstep{按份数公式求边}{$PA=PC\times\dfrac{3}{4}=8\times\dfrac{3}{4}=6$。}
\end{solutionblock}
\subsection*{第 4 题}
\setcounter{solstep}{0}
\begin{solutionblock}
\textbf{答案：}$CD=10$。\par\medskip
    \solstep{先判断}{已知边与未知边中出现平行线段 $AB$、$CD$，先求两个相似三角形的比。}
    \solstep{计算并约分}{$\dfrac{BP}{PD}=\dfrac{6}{15}=\dfrac{2}{5}$。}
    \solstep{标份数}{约分后，$CD$ 对应 5 份，$AB$ 对应 2 份。}
    \solstep{按份数公式求边}{$CD=AB\times\dfrac{5}{2}=4\times\dfrac{5}{2}=10$。}
\end{solutionblock}
\subsection*{第 5 题}
\setcounter{solstep}{0}
\begin{solutionblock}
\textbf{答案：}$CD=25$。\par\medskip
    \solstep{先判断}{已知边与未知边中出现平行线段 $AB$、$CD$，先求两个相似三角形的比。}
    \solstep{计算并约分}{$\dfrac{PA}{PC}=\dfrac{6}{10}=\dfrac{3}{5}$。}
    \solstep{标份数}{约分后，$CD$ 对应 5 份，$AB$ 对应 3 份。}
    \solstep{按份数公式求边}{$CD=AB\times\dfrac{5}{3}=15\times\dfrac{5}{3}=25$。}
\end{solutionblock}
\subsection*{第 6 题}
\setcounter{solstep}{0}
\begin{solutionblock}
\textbf{答案：}$PD=10$。\par\medskip
    \solstep{先判断}{已知边与未知边都不在平行线段 $AB,CD$ 上，求两侧对应线段的比。}
    \solstep{计算并约分}{$\dfrac{AP}{PC}=\dfrac{8}{10}=\dfrac{4}{5}$。}
    \solstep{标份数}{约分后，$PD$ 对应 5 份，$BP$ 对应 4 份。}
    \solstep{按份数公式求边}{$PD=BP\times\dfrac{5}{4}=8\times\dfrac{5}{4}=10$。}
\end{solutionblock}
\subsection*{第 7 题}
\setcounter{solstep}{0}
\begin{solutionblock}
\textbf{答案：}$PC=14$。\par\medskip
    \solstep{先判断}{已知边与未知边都不在平行线段 $AB,CD$ 上，求两侧对应线段的比。}
    \solstep{计算并约分}{$\dfrac{PB}{PD}=\dfrac{6}{14}=\dfrac{3}{7}$。}
    \solstep{标份数}{约分后，$PC$ 对应 7 份，$PA$ 对应 3 份。}
    \solstep{按份数公式求边}{$PC=PA\times\dfrac{7}{3}=6\times\dfrac{7}{3}=14$。}
\end{solutionblock}
\subsection*{第 8 题}
\setcounter{solstep}{0}
\begin{solutionblock}
\textbf{答案：}$AB=16$。\par\medskip
    \solstep{先判断}{已知边与未知边中出现平行线段 $CD$、$AB$，先求两个相似三角形的比。}
    \solstep{计算并约分}{$\dfrac{BP}{PD}=\dfrac{8}{14}=\dfrac{4}{7}$。}
    \solstep{标份数}{约分后，$AB$ 对应 4 份，$CD$ 对应 7 份。}
    \solstep{按份数公式求边}{$AB=CD\times\dfrac{4}{7}=28\times\dfrac{4}{7}=16$。}
\end{solutionblock}
\subsection*{第 9 题}
\setcounter{solstep}{0}
\begin{solutionblock}
\textbf{答案：}$AB=15$。\par\medskip
    \solstep{先判断}{已知边与未知边中出现平行线段 $CD$、$AB$，先求两个相似三角形的比。}
    \solstep{计算并约分}{$\dfrac{PA}{PC}=\dfrac{10}{16}=\dfrac{5}{8}$。}
    \solstep{标份数}{约分后，$AB$ 对应 5 份，$CD$ 对应 8 份。}
    \solstep{按份数公式求边}{$AB=CD\times\dfrac{5}{8}=24\times\dfrac{5}{8}=15$。}
\end{solutionblock}
\subsection*{第 10 题}
\setcounter{solstep}{0}
\begin{solutionblock}
\textbf{答案：}$BP=10$。\par\medskip
    \solstep{先判断}{已知边与未知边都不在平行线段 $AB,CD$ 上，求两侧对应线段的比。}
    \solstep{计算并约分}{$\dfrac{AP}{PC}=\dfrac{10}{18}=\dfrac{5}{9}$。}
    \solstep{标份数}{约分后，$BP$ 对应 5 份，$PD$ 对应 9 份。}
    \solstep{按份数公式求边}{$BP=PD\times\dfrac{5}{9}=18\times\dfrac{5}{9}=10$。}
\end{solutionblock}
\clearpage
\section*{练习 2 答案}
\subsection*{第 1 题}
\setcounter{solstep}{0}
\begin{solutionblock}
\textbf{答案：}$AB=\frac{12}{7}$。\par\medskip
    \solstep{先判断}{已知边与未知边中出现平行线段 $CD$、$AB$，先求两个相似三角形的比。}
    \solstep{计算并约分}{$\dfrac{PA}{PC}=\dfrac{\frac{6}{7}}{\frac{36}{7}}=\dfrac{1}{6}$。}
    \solstep{标份数}{约分后，$AB$ 对应 $1$ 份，$CD$ 对应 $6$ 份。}
    \solstep{按份数公式求边}{$AB=CD\times\dfrac{1}{6}=\frac{72}{7}\times\dfrac{1}{6}=\frac{12}{7}$。}
\end{solutionblock}
\subsection*{第 2 题}
\setcounter{solstep}{0}
\begin{solutionblock}
\textbf{答案：}$BP=3\sqrt{2}$。\par\medskip
    \solstep{先判断}{已知边与未知边都不在平行线段 $AB,CD$ 上，先求两侧对应线段的比。}
    \solstep{计算并约分}{$\dfrac{AP}{PC}=\dfrac{\sqrt{2}}{2\sqrt{2}}=\dfrac{1}{2}$。}
    \solstep{标份数}{约分后，$BP$ 对应 $1$ 份，$PD$ 对应 $2$ 份。}
    \solstep{按份数公式求边}{$BP=PD\times\dfrac{1}{2}=6\sqrt{2}\times\dfrac{1}{2}=3\sqrt{2}$。}
\end{solutionblock}
\subsection*{第 3 题}
\setcounter{solstep}{0}
\begin{solutionblock}
\textbf{答案：}$PA=2\sqrt{2}$。\par\medskip
    \solstep{先判断}{已知边与未知边都不在平行线段 $AB,CD$ 上，先求两侧对应线段的比。}
    \solstep{计算并约分}{$\dfrac{PB}{PD}=\dfrac{\sqrt{2}}{\sqrt{14}}=\dfrac{1}{\sqrt{7}}$。}
    \solstep{标份数}{约分后，$PA$ 对应 $1$ 份，$PC$ 对应 $\sqrt{7}$ 份。}
    \solstep{按份数公式求边}{$PA=PC\times\dfrac{1}{\sqrt{7}}=2\sqrt{14}\times\dfrac{1}{\sqrt{7}}=2\sqrt{2}$。}
\end{solutionblock}
\subsection*{第 4 题}
\setcounter{solstep}{0}
\begin{solutionblock}
\textbf{答案：}$CD=\frac{3}{2}$。\par\medskip
    \solstep{先判断}{已知边与未知边中出现平行线段 $AB$、$CD$，先求两个相似三角形的比。}
    \solstep{计算并约分}{$\dfrac{BP}{PD}=\dfrac{\frac{1}{8}}{\frac{3}{8}}=\dfrac{1}{3}$。}
    \solstep{标份数}{约分后，$CD$ 对应 $3$ 份，$AB$ 对应 $1$ 份。}
    \solstep{按份数公式求边}{$CD=AB\times\dfrac{3}{1}=\frac{1}{2}\times\dfrac{3}{1}=\frac{3}{2}$。}
\end{solutionblock}
\subsection*{第 5 题}
\setcounter{solstep}{0}
\begin{solutionblock}
\textbf{答案：}$CD=\frac{63}{4}$。\par\medskip
    \solstep{先判断}{已知边与未知边中出现平行线段 $AB$、$CD$，先求两个相似三角形的比。}
    \solstep{计算并约分}{$\dfrac{PA}{PC}=\dfrac{\frac{7}{8}}{\frac{21}{4}}=\dfrac{1}{6}$。}
    \solstep{标份数}{约分后，$CD$ 对应 $6$ 份，$AB$ 对应 $1$ 份。}
    \solstep{按份数公式求边}{$CD=AB\times\dfrac{6}{1}=\frac{21}{8}\times\dfrac{6}{1}=\frac{63}{4}$。}
\end{solutionblock}
\subsection*{第 6 题}
\setcounter{solstep}{0}
\begin{solutionblock}
\textbf{答案：}$PD=4\sqrt{15}$。\par\medskip
    \solstep{先判断}{已知边与未知边都不在平行线段 $AB,CD$ 上，先求两侧对应线段的比。}
    \solstep{计算并约分}{$\dfrac{AP}{PC}=\dfrac{\sqrt{10}}{2\sqrt{15}}=\dfrac{1}{\sqrt{6}}$。}
    \solstep{标份数}{约分后，$PD$ 对应 $\sqrt{6}$ 份，$BP$ 对应 $1$ 份。}
    \solstep{按份数公式求边}{$PD=BP\times\dfrac{\sqrt{6}}{1}=2\sqrt{10}\times\dfrac{\sqrt{6}}{1}=4\sqrt{15}$。}
\end{solutionblock}
\subsection*{第 7 题}
\setcounter{solstep}{0}
\begin{solutionblock}
\textbf{答案：}$PC=3\sqrt{30}$。\par\medskip
    \solstep{先判断}{已知边与未知边都不在平行线段 $AB,CD$ 上，先求两侧对应线段的比。}
    \solstep{计算并约分}{$\dfrac{PB}{PD}=\dfrac{\sqrt{5}}{\sqrt{30}}=\dfrac{1}{\sqrt{6}}$。}
    \solstep{标份数}{约分后，$PC$ 对应 $\sqrt{6}$ 份，$PA$ 对应 $1$ 份。}
    \solstep{按份数公式求边}{$PC=PA\times\dfrac{\sqrt{6}}{1}=3\sqrt{5}\times\dfrac{\sqrt{6}}{1}=3\sqrt{30}$。}
\end{solutionblock}
\subsection*{第 8 题}
\setcounter{solstep}{0}
\begin{solutionblock}
\textbf{答案：}$AB=5$。\par\medskip
    \solstep{先判断}{已知边与未知边中出现平行线段 $CD$、$AB$，先求两个相似三角形的比。}
    \solstep{计算并约分}{$\dfrac{BP}{PD}=\dfrac{\frac{5}{4}}{5}=\dfrac{1}{4}$。}
    \solstep{标份数}{约分后，$AB$ 对应 $1$ 份，$CD$ 对应 $4$ 份。}
    \solstep{按份数公式求边}{$AB=CD\times\dfrac{1}{4}=20\times\dfrac{1}{4}=5$。}
\end{solutionblock}
\subsection*{第 9 题}
\setcounter{solstep}{0}
\begin{solutionblock}
\textbf{答案：}$AB=\frac{7}{3}$。\par\medskip
    \solstep{先判断}{已知边与未知边中出现平行线段 $CD$、$AB$，先求两个相似三角形的比。}
    \solstep{计算并约分}{$\dfrac{PA}{PC}=\dfrac{\frac{7}{6}}{\frac{14}{3}}=\dfrac{1}{4}$。}
    \solstep{标份数}{约分后，$AB$ 对应 $1$ 份，$CD$ 对应 $4$ 份。}
    \solstep{按份数公式求边}{$AB=CD\times\dfrac{1}{4}=\frac{28}{3}\times\dfrac{1}{4}=\frac{7}{3}$。}
\end{solutionblock}
\subsection*{第 10 题}
\setcounter{solstep}{0}
\begin{solutionblock}
\textbf{答案：}$BP=3\sqrt{3}$。\par\medskip
    \solstep{先判断}{已知边与未知边都不在平行线段 $AB,CD$ 上，先求两侧对应线段的比。}
    \solstep{计算并约分}{$\dfrac{AP}{PC}=\dfrac{\sqrt{3}}{3\sqrt{2}}=\dfrac{1}{\sqrt{6}}$。}
    \solstep{标份数}{约分后，$BP$ 对应 $1$ 份，$PD$ 对应 $\sqrt{6}$ 份。}
    \solstep{按份数公式求边}{$BP=PD\times\dfrac{1}{\sqrt{6}}=9\sqrt{2}\times\dfrac{1}{\sqrt{6}}=3\sqrt{3}$。}
\end{solutionblock}
\clearpage
\section*{练习 3 答案}
\subsection*{第 1 题}
\setcounter{solstep}{0}
\begin{solutionblock}
\textbf{答案：}$PA=\frac{3}{8}$。\par\medskip
    \solstep{先判断}{已知边与未知边都不在平行线段 $AB,CD$ 上，先求两侧对应线段的比。}
    \solstep{计算并约分}{$\dfrac{PB}{PD}=\dfrac{\frac{1}{8}}{\frac{5}{8}}=\dfrac{1}{5}$。}
    \solstep{标份数}{约分后，$PA$ 对应 $1$ 份，$PC$ 对应 $5$ 份。}
    \solstep{按份数公式求边}{$PA=PC\times\dfrac{1}{5}=\frac{15}{8}\times\dfrac{1}{5}=\frac{3}{8}$。}
\end{solutionblock}
\subsection*{第 2 题}
\setcounter{solstep}{0}
\begin{solutionblock}
\textbf{答案：}$CD=14$。\par\medskip
    \solstep{先判断}{已知边与未知边中出现平行线段 $AB$、$CD$，先求两个相似三角形的比。}
    \solstep{计算并约分}{$\dfrac{BP}{PD}=\dfrac{\sqrt{7}}{7}=\dfrac{1}{\sqrt{7}}$。}
    \solstep{标份数}{约分后，$CD$ 对应 $\sqrt{7}$ 份，$AB$ 对应 $1$ 份。}
    \solstep{按份数公式求边}{$CD=AB\times\dfrac{\sqrt{7}}{1}=2\sqrt{7}\times\dfrac{\sqrt{7}}{1}=14$。}
\end{solutionblock}
\subsection*{第 3 题}
\setcounter{solstep}{0}
\begin{solutionblock}
\textbf{答案：}$CD=24$。\par\medskip
    \solstep{先判断}{已知边与未知边中出现平行线段 $AB$、$CD$，先求两个相似三角形的比。}
    \solstep{计算并约分}{$\dfrac{PA}{PC}=\dfrac{\sqrt{6}}{6}=\dfrac{1}{\sqrt{6}}$。}
    \solstep{标份数}{约分后，$CD$ 对应 $\sqrt{6}$ 份，$AB$ 对应 $1$ 份。}
    \solstep{按份数公式求边}{$CD=AB\times\dfrac{\sqrt{6}}{1}=4\sqrt{6}\times\dfrac{\sqrt{6}}{1}=24$。}
\end{solutionblock}
\subsection*{第 4 题}
\setcounter{solstep}{0}
\begin{solutionblock}
\textbf{答案：}$PD=21$。\par\medskip
    \solstep{先判断}{已知边与未知边都不在平行线段 $AB,CD$ 上，先求两侧对应线段的比。}
    \solstep{计算并约分}{$\dfrac{AP}{PC}=\dfrac{\frac{7}{4}}{\frac{21}{2}}=\dfrac{1}{6}$。}
    \solstep{标份数}{约分后，$PD$ 对应 $6$ 份，$BP$ 对应 $1$ 份。}
    \solstep{按份数公式求边}{$PD=BP\times\dfrac{6}{1}=\frac{7}{2}\times\dfrac{6}{1}=21$。}
\end{solutionblock}
\subsection*{第 5 题}
\setcounter{solstep}{0}
\begin{solutionblock}
\textbf{答案：}$PC=\frac{6}{7}$。\par\medskip
    \solstep{先判断}{已知边与未知边都不在平行线段 $AB,CD$ 上，先求两侧对应线段的比。}
    \solstep{计算并约分}{$\dfrac{PB}{PD}=\dfrac{\frac{1}{7}}{\frac{2}{7}}=\dfrac{1}{2}$。}
    \solstep{标份数}{约分后，$PC$ 对应 $2$ 份，$PA$ 对应 $1$ 份。}
    \solstep{按份数公式求边}{$PC=PA\times\dfrac{2}{1}=\frac{3}{7}\times\dfrac{2}{1}=\frac{6}{7}$。}
\end{solutionblock}
\subsection*{第 6 题}
\setcounter{solstep}{0}
\begin{solutionblock}
\textbf{答案：}$AB=4\sqrt{3}$。\par\medskip
    \solstep{先判断}{已知边与未知边中出现平行线段 $CD$、$AB$，先求两个相似三角形的比。}
    \solstep{计算并约分}{$\dfrac{BP}{PD}=\dfrac{\sqrt{3}}{3}=\dfrac{1}{\sqrt{3}}$。}
    \solstep{标份数}{约分后，$AB$ 对应 $1$ 份，$CD$ 对应 $\sqrt{3}$ 份。}
    \solstep{按份数公式求边}{$AB=CD\times\dfrac{1}{\sqrt{3}}=12\times\dfrac{1}{\sqrt{3}}=4\sqrt{3}$。}
\end{solutionblock}
\subsection*{第 7 题}
\setcounter{solstep}{0}
\begin{solutionblock}
\textbf{答案：}$AB=4\sqrt{2}$。\par\medskip
    \solstep{先判断}{已知边与未知边中出现平行线段 $CD$、$AB$，先求两个相似三角形的比。}
    \solstep{计算并约分}{$\dfrac{PA}{PC}=\dfrac{2\sqrt{2}}{8}=\dfrac{1}{2 \sqrt{2}}$。}
    \solstep{标份数}{约分后，$AB$ 对应 $1$ 份，$CD$ 对应 $2 \sqrt{2}$ 份。}
    \solstep{按份数公式求边}{$AB=CD\times\dfrac{1}{2 \sqrt{2}}=16\times\dfrac{1}{2 \sqrt{2}}=4\sqrt{2}$。}
\end{solutionblock}
\subsection*{第 8 题}
\setcounter{solstep}{0}
\begin{solutionblock}
\textbf{答案：}$BP=\frac{9}{7}$。\par\medskip
    \solstep{先判断}{已知边与未知边都不在平行线段 $AB,CD$ 上，先求两侧对应线段的比。}
    \solstep{计算并约分}{$\dfrac{AP}{PC}=\dfrac{\frac{3}{7}}{\frac{12}{7}}=\dfrac{1}{4}$。}
    \solstep{标份数}{约分后，$BP$ 对应 $1$ 份，$PD$ 对应 $4$ 份。}
    \solstep{按份数公式求边}{$BP=PD\times\dfrac{1}{4}=\frac{36}{7}\times\dfrac{1}{4}=\frac{9}{7}$。}
\end{solutionblock}
\subsection*{第 9 题}
\setcounter{solstep}{0}
\begin{solutionblock}
\textbf{答案：}$PA=\frac{36}{5}$。\par\medskip
    \solstep{先判断}{已知边与未知边都不在平行线段 $AB,CD$ 上，先求两侧对应线段的比。}
    \solstep{计算并约分}{$\dfrac{PB}{PD}=\dfrac{\frac{9}{5}}{9}=\dfrac{1}{5}$。}
    \solstep{标份数}{约分后，$PA$ 对应 $1$ 份，$PC$ 对应 $5$ 份。}
    \solstep{按份数公式求边}{$PA=PC\times\dfrac{1}{5}=36\times\dfrac{1}{5}=\frac{36}{5}$。}
\end{solutionblock}
\subsection*{第 10 题}
\setcounter{solstep}{0}
\begin{solutionblock}
\textbf{答案：}$CD=4\sqrt{10}$。\par\medskip
    \solstep{先判断}{已知边与未知边中出现平行线段 $AB$、$CD$，先求两个相似三角形的比。}
    \solstep{计算并约分}{$\dfrac{BP}{PD}=\dfrac{\sqrt{10}}{2\sqrt{10}}=\dfrac{1}{2}$。}
    \solstep{标份数}{约分后，$CD$ 对应 $2$ 份，$AB$ 对应 $1$ 份。}
    \solstep{按份数公式求边}{$CD=AB\times\dfrac{2}{1}=2\sqrt{10}\times\dfrac{2}{1}=4\sqrt{10}$。}
\end{solutionblock}
\end{document}
